\begin{table}[!tbp]\centering
\caption{Estadísticas descriptivas por grupo educativo, período previo a la reforma}\label{tab:desc}
\begin{threeparttable}
\begin{tabular}{lccc}
\toprule
 & Baja educación & Alta educación & Diferencia \\
 & ($\le$ sec. completa) & ($\ge$ superior) & \\
\midrule
Edad & 35.92 & 34.65 & 1.27 \\
Mujer & 0.50 & 0.50 & 0.00 \\
Urbano & 0.76 & 0.93 & -0.17 \\
Casado(a) & 0.53 & 0.41 & 0.12 \\
Años de educación & 8.57 & 14.47 & -5.90 \\
Participación laboral & 0.73 & 0.80 & -0.07 \\
Ocupado & 0.69 & 0.73 & -0.04 \\
Asalariado (incl. del hogar) & 0.29 & 0.48 & -0.19 \\
Independiente & 0.39 & 0.24 & 0.15 \\
Formal (si está ocupado) & 0.13 & 0.47 & -0.34 \\
Informal (si está ocupado) & 0.87 & 0.53 & 0.34 \\
Horas semanales (si está ocupado) & 39.74 & 41.44 & -1.70 \\
Salario real por hora (S/) & 5.74 & 9.08 & -3.34 \\
Ingreso laboral real mensual (S/) & 904.30 & 1\,755.12 & -850.82 \\
Bajo el SM (asalariados) & 0.43 & 0.19 & 0.24 \\
\midrule
Observaciones & 65\,954 & 27\,278 & \\
\bottomrule
\end{tabular}
\begin{tablenotes}\footnotesize
\item Notas: Medias ponderadas con los factores de expansión de la ENAHO; personas en edad de trabajar (14--65) en los trimestres previos a la reforma (2021Q1--2022Q1). Baja educación: a lo sumo secundaria completa (\texttt{p301a}$\le$6). Alta educación: cualquier nivel de educación superior (\texttt{p301a}$\ge$7). Las variables de asalariados (salario por hora, bajo el SM) se condicionan a ser asalariado del sector privado (empleados y trabajadores del hogar, ingresos mensuales equivalentes); las horas y la formalidad se condicionan a estar ocupado. Los conteos de observaciones no están ponderados. La fila de bajo el SM usa el mínimo legal vigente en el mes de referencia del ingreso (930 soles durante toda la ventana previa a la reforma); las proporciones de exposición respecto del NUEVO mínimo (57 frente a 40 por ciento) que se discuten en el texto usan 1025 soles.
\end{tablenotes}
\end{threeparttable}\end{table}
